\documentclass{article}
\usepackage{graphicx} % Required for inserting images
% Pacotes matemáticos
\usepackage{amsmath, amssymb, amsthm}
\usepackage[a4paper,margin=2.5cm]{geometry}

% Ambientes de teorema
\newtheorem{theorem}{Teorema}[section]
\newtheorem{lemma}[theorem]{Lema}
\newtheorem{proposition}[theorem]{Proposição}
\newtheorem{corollary}[theorem]{Corolário}

\theoremstyle{definition}
\newtheorem{definition}[theorem]{Definição}
\newtheorem{example}[theorem]{Exemplo}

\theoremstyle{remark}
\newtheorem{remark}[theorem]{Observação}
\title{Completude de $S_n$}
\author{Gustavo Reis de Araújo Pinheiro\\Andrés David Cadena Simons\\
Jean da Cruz Oliveira\\
\large Universidade Federal de Minais Gerais}

\date{Junho 2026}

\begin{document}

\maketitle
    Grupos completos ocupam um lugar de destaque na teoria dos grupos, pois coincidem com o seu próprio grupo de automorfismos. Nesta palestra estudaremos a completude dos grupos de permutação, demonstrando que $S_n$ é completo para $n\neq 6$. Para isso, apresentaremos um lema técnico sobre os automorfismos de $S_n$ e analisaremos como esse resultado se relaciona com o problema da torre de automorfismos. Por fim, discutiremos o caráter excepcional de $S_6$, cuja estrutura de automorfismos difere daquela dos demais grupos de permutação.
\section{Motivação: O problema da torre dos automorfismos.}
Dado um grupo $G$, pelo teorema do isomorfismo de grupos; 
  \begin{align*}
      G/Z(G)\cong Inn(G) \lhd Aut(G).
    \end{align*}
    \begin{lemma}
         Dado $G$ grupo, temos que se $Z(G)=\{1\}$, então $Z(Aut(G))=\{1\}$.
    \end{lemma}
   Assim, seja $G$ grupo com $Z(G)=\{1\}$, denotando $G_0=G$. Pelo lema anterior temos que $G_0\lhd Aut(G_0)$. Depois, denotando $G_1=Aut(G_0)$, temos que como $Z(G_1)=\{1\}$, então $G_0\lhd G_1\lhd Aut(G_1)$. Assim, raciocinando de maneira indutiva temos
    \begin{align*}
      G_0\lhd \underbrace{G_1}_{Aut(G_0)}\lhd \underbrace{G_2}_{Aut(G_1)}\lhd\underbrace{G_3}_{Aut(G_2)}\lhd\cdots\lhd G_{\alpha}\lhd G_{\alpha+1}\lhd\cdots
    \end{align*}
    com $\alpha$ ordinal.  Esta cadeia ascendente é chamada a \textbf{torre de automorfismos}.

    Pelo Teorema de Wielandt, a torre de automorfismos de $G$ estabiliza num número finito de passos. 

    \begin{definition}
        Um grupo $G$ é dito \textbf{completo} quando $Z(G) = 1$ e $Aut(G) = Inn(G)$
    \end{definition}
Note que $G_{N}=G_{N+1}$ é o mesmo que dizer que $Inn(G_N)=Aut(G_{N})$, pelo que podemos assegurar que a cadeia estabiliza num grupo \textbf{completo}.

\section{Completude de $S_n$}
\begin{lemma}
      Sejam $n\neq 6$ e $\tau \in S_n$ um elemento de ordem $2$ que é produto de $k$ transposições disjuntas. Então:

    \begin{enumerate}
        \item O tamanho da classe de conjugação de $\tau$ é
        $$\frac{n!}{2^k(n-2k)!k!}$$

        \item Se $k>1$ então $\vert Cl(\tau)\vert \neq \vert Cl(\sigma)\vert$, para qualquer transposição $\sigma \in S_n$.
\end{enumerate}
\end{lemma}
   
\begin{proposition}
     Para $n\geq 3,$ temos que $Z(S_n)=\{1\}.$
\end{proposition}

\begin{theorem}
        Se $n \neq 6$, todo automorfismo de $S_n$ é interno.
\end{theorem}
\section{$S_6$ não é completo}
\begin{definition}
  Dado um conjunto $X$, dizemos que um subgrupo $S$ de $Sim(X)$ é um \textbf{subgrupo transitivo} se dados $x,y\in X$, temos que existe $\sigma\in S$ tal que $\sigma(x)=y$.     
\end{definition}
\begin{lemma}[]
  Seja $G$ um grupo e $H\leq G$. Definamos $X=\{xH:x\in G\}$ (o conjunto das classes laterais a esquerda de $H$) e $\varphi$ da seguinte maneira
  \begin{align*}
    \varphi:G&\to Sim(X),\\
            g&\to \varphi_g:X\to X,\\
  \phantom{g}&\phantom{\to \varphi_g:}xH\to gxH.
  \end{align*}
  Então $\Img \varphi$ é um subgrupo transitivo de $Sim(X)$.
\end{lemma}
\begin{lemma}
      Existe um subgrupo transitivo $K$ de $S_6$ com ordem $120$ que não contém transposições. 
\end{lemma}
\begin{theorem}
       Existe um automorfismo de $S_6$ que não é interno.     
\end{theorem}

\end{document}
